\documentclass[12pt]{article}
\usepackage[utf8]{inputenc}
\usepackage[T1]{fontenc}
\usepackage[margin=1in]{geometry}
\usepackage{setspace}
\usepackage{xcolor}
\usepackage[colorlinks=true, urlcolor=blue!60!black, linkcolor=blue!60!black]{hyperref}
\setlength{\parindent}{2em}
\setlength{\parskip}{1.25em}
\renewcommand{\baselinestretch}{1.0}
\begin{document}
\begin{center}
{\LARGE \textit{Gonzales v. Raich}, 545 U.S. 1 (2005)\par}
\vspace{0.8em}
{\large Case Brief --- Four-Page Edition\par}
\end{center}
\vspace{1em}

\section{Memory Jogger}

Under the Commerce and Necessary and Proper Clauses, Congress may criminalize purely intrastate, noncommercial cultivation and possession of homegrown medical marijuana authorized by state law, because regulating local production of a fungible commodity is rationally related to a comprehensive federal regime controlling the interstate drug market (\href{https://www.courtlistener.com/opinion/799995/gonzales-v-raich/}{CourtListener}).

\section{Detailed Case Facts}

California's Compassionate Use Act of 1996 (Proposition 215) exempted seriously ill patients and their caregivers from state prosecution for possessing or cultivating marijuana on a physician's recommendation. Angel Raich and Diane Monson suffered serious medical conditions; their physicians concluded that cannabis was the only effective treatment, and Raich's physician attested that forgoing it could prove fatal. Monson grew her own plants; Raich received locally grown cannabis free from two caregivers. The federal Controlled Substances Act (CSA) classifies marijuana as a Schedule I substance and makes its manufacture, distribution, or possession a crime. In August 2002, although county officials found that Monson's use complied with California law, DEA agents seized and destroyed her six plants.

\section{Procedural History}

Raich, Monson, and the caregivers sued the Attorney General for declaratory and injunctive relief. The district court denied a preliminary injunction. A divided Ninth Circuit reversed, treating ``the intrastate, noncommercial cultivation and possession of cannabis for personal medical purposes'' as a separate class beyond the commerce power under \textit{Lopez} and \textit{Morrison}. The Supreme Court vacated and remanded on June 6, 2005, leaving the due process and medical-necessity claims open. Judge Beam, dissenting below, had found the case indistinguishable from \textit{Wickard}.

\section{Judicial Votes}

Vacated and remanded, \textbf{6--3}.
\begin{itemize}
\item \textbf{Majority:} Justice Stevens, joined by Justices Kennedy, Souter, Ginsburg, and Breyer.
\item \textbf{Concurring in the judgment:} Justice Scalia.
\item \textbf{Dissents:} Justice O'Connor, joined by Chief Justice Rehnquist and Justice Thomas (except Part III); Justice Thomas separately.
\end{itemize}

\section{Holding}

The CSA's ban on manufacturing and possessing marijuana is constitutional as applied to intrastate, noncommercial medical cultivation and use lawful under California law. Congress may regulate a purely local class of activities involving a fungible commodity when it has a rational basis to conclude that leaving it unregulated would undercut a comprehensive scheme governing the interstate market, and courts may not excise individual applications from such a scheme.

\section{Analysis of Opinions}

\textbf{Majority (Stevens).} The case falls in the third \textit{Lopez} category, activities substantially affecting interstate commerce, and parallels \textit{Wickard v. Filburn}: like home-consumed wheat, homegrown marijuana is a fungible commodity with an interstate market. Under rational-basis review, Congress could conclude that exempting local cultivation, hard to distinguish from interstate supply and prone to diversion, would leave a ``gaping hole'' in the CSA. \textit{Lopez} and \textit{Morrison} were facial challenges to single-subject statutes regulating noneconomic conduct; here respondents sought to excise a slice of a valid comprehensive scheme. Defining ``economics'' as ``the production, distribution, and consumption of commodities,'' Stevens found the CSA quintessentially economic. Under the Supremacy Clause, state law can neither expand nor narrow Congress's power. In dicta, Stevens suggested that respondents' medical evidence could cast doubt on the Schedule I finding and pointed them to administrative rescheduling or Congress.

\textbf{Scalia (concurring in the judgment).} The power to reach local activity comes from the Necessary and Proper Clause. Under \textit{McCulloch} and \textit{Darby}, Congress may regulate even noneconomic local activity when doing so is a ``plainly adapted'' means of making interstate regulation effective; marijuana is ``never more than an instant from the interstate market.''

\textbf{O'Connor (dissenting).} Invoking States as ``laboratories'' of democracy, she warned that the comprehensive-scheme rationale makes \textit{Lopez} a mere drafting guide, called the majority's definition of ``economic'' ``breathtaking,'' and noted the absence of evidence that California's medical users affected the national market.

\textbf{Thomas (dissenting).} ``Commerce'' originally meant trade, not production or home cultivation. Applying the CSA here was neither ``necessary'' nor ``proper,'' and under the majority's logic, ``the Federal Government may now regulate quilting bees, clothes drives, and potluck suppers throughout the 50 States.''

\section{Examples: Future Applications}

\begin{enumerate}
\item \textbf{Same-Side (Upheld): Homemade Machine Guns.} A comprehensive firearms statute bans home assembly and possession of machine guns. Firearms are fungible commodities with an interstate market, so \textit{Raich} sustains the ban.
\item \textbf{Same-Side (Upheld): State-Legal Psilocybin.} A resident grows psilocybin in his basement under a state program. The CSA reaches it regardless of state authorization.
\item \textbf{Opposite-Side (Struck Down): Backyard Composting.} A single-subject statute bans composting near schools, with no broader scheme or jurisdictional element. Like \textit{Lopez}, it targets noncommercial local conduct.
\item \textbf{Opposite-Side (Struck Down): Broccoli Mandate.} A law compels adults to buy broccoli. Under \textit{NFIB v. Sebelius} (2012), \textit{Raich} reaches existing economic activity, not inactivity.
\item \textbf{Fence-Sitter: A One-County Cave Beetle.} The Endangered Species Act bars clearing brush to protect a beetle never bought or sold. The government invokes the comprehensive scheme, but the beetle is no fungible commodity.
\end{enumerate}

\section{Critique}

\textbf{Originalist critique.} Randy Barnett, who argued for respondents, contended that \textit{Raich} abandoned the founding meaning of ``commerce'' as trade, though Scalia's concurrence preserved some limits on pretextual federal police power. The Court later cabined \textit{Raich} in \textit{NFIB v. Sebelius}.

\textbf{Doctrinal and progressive critique.} \textit{Raich}'s irony was a liberal bloc defending New Deal power to sustain the War on Drugs against progressive state reforms. Commentators say it left doctrine muddled, since local regulations survive because they sit inside larger schemes, and it led some progressives to value federalism as a shield for state innovation.

\section{Key Quotations}

\begin{quote}
``We need not determine whether respondents' activities, taken in the aggregate, substantially affect interstate commerce in fact, but only whether a `rational basis' exists for so concluding.''
\end{quote}
\noindent --- Justice Stevens, for the Court, 545 U.S. 1.

\begin{quote}
``Where necessary to make a regulation of interstate commerce effective, Congress may regulate even those intrastate activities that do not themselves substantially affect interstate commerce.''
\end{quote}
\noindent --- Justice Scalia, concurring in the judgment, 545 U.S. 1.

\begin{quote}
``The similarities between this case and \textit{Wickard} are striking. Like the farmer in \textit{Wickard}, respondents are cultivating, for home consumption, a fungible commodity for which there is an established, albeit illegal, interstate market.''
\end{quote}
\noindent --- Justice Stevens, for the Court, 545 U.S. 1.

\end{document}
